% Chapter 2

\chapter{Estado da Arte} % Main chapter title

\label{chap:Chapter2} % For referencing the chapter elsewhere, use \ref{chap:Chapter2} 

Este capítulo apresenta uma revisão da literatura no domínio dos sistemas gimbal para seguimento de UAVs. A análise evolui de uma perspetiva macro sobre sistemas anti-drones com a utilização gimbal, detalhando em seguida a deteção e o seguimento. Esta abordagem visa otimizar a performance da prova de conceito através da integração de mecanismos mais recentes na literatura.


\section{Sistemas Gimbal Anti-Drone}

Devido à elevada especificidade e ao carácter estratégico dos sistemas de Energia Direcionada, a literatura técnica detalhada sobre sistemas de intersecção laser para fins de defesa é restrita. Todavia, o setor industrial de defesa já apresenta soluções consolidadas desenvolvidas por empresas de referência. 

No âmbito académico, destaca-se o trabalho desenvolvido por \cite{sistemalasergimbal}, que descreve um sistema de seguimento em dois eixos fundamentado em visão estereoscópica através de uma câmara de profundidade Intel RealSense D435. O sistema integra o algoritmo YOLOv5 para a deteção e localização do alvo em tempo real. Em termos de desempenho, o estudo reportou limitações na confiança do modelo de deteção proporcionais à distância, registando valores inferiores a 20\% para drones de pequeno porte (23 cm de envergadura) a 8 m, em contraste com os cerca de 55\% obtidos para drones de maior porte (69 cm de envergadura) à mesma distância. No que concerne à exatidão dinâmica, o sistema revelou-se capaz de acompanhar alvos com velocidades entre 2 m/s e 4 m/s, apresentando um erro de desvio entre 1,9 cm a 4 m e 6,2 cm a uma distância de 10 m, o que valida a eficácia da intersecção laser mesmo perante a mobilidade moderada do alvo.

Outro estudo que se enquadra neste âmbito, embora não utilize laser, é o trabalho de \cite{camera_gimbal}. Os autores apresentam um sistema gimbal para deteção e seguimento de drones baseado no modelo personalizado Drone-Net, que utiliza o framework Faster R-CNN (um algoritmo de deteção de dois estágios). Para o seguimento, o artigo propõe a substituição do controlo PID por um algoritmo de controlo adaptativo de velocidade. Este método calcula o erro de píxeis entre o centro da imagem e a localização do alvo, aplicando uma de três velocidades predefinidas para corrigir a postura do gimbal de forma mais estável e rápida. Diferente do trabalho anterior, este estudo não verifica a distância máxima de deteção e seguimento, mas verifica a robustez do sistema em distinguir drones de interferências comuns, como pássaros e papagaios de papel, utilizando o conjunto de dados \emph{Drone-Kite-Bird}. Em termos de resultados, o modelo Drone-Net alcançou um média da exatidão (mAP) de 69.1\%. No que respeita à exatidão de seguimento, o sistema demonstrou ser capaz de manter o drone posicionado na área central de 5\% do campo de visão em 71.2\% dos quadros analisados no cenário de teste a 80 metros, comprovando a eficácia do controlo adaptativo em trajetórias de longa distância.

O mercado de defesa demonstra o interesse crescente desta solução. Diversas empresas globais de referência já implementaram ou encontram-se em fases avançadas de desenvolvimento de produtos que utilizam lasers de alta energia para a intersecção de ameaças aéreas. A empresa \cite{lockheed2021helios}, com o seu sistema HELIOS, e a \cite{rafael2023navalironbeam}, com o seu sistema Naval Iron Beam, permitem a integração em plataformas navais, utilizando sistemas gimbal de alta exatidão acoplados a emissores laser e sensores de deteção próprios. Por outro lado, empresas como a \cite{rtx2023hel} apresentam sistemas portáteis com capacidades para detetar e neutralizar UAVs em diversos cenários. Todavia, nenhuma destas empresas disponibiliza informação técnica detalhada que possa ser utilizada diretamente para o desenvolvimento de investigação aberta, o que reforça a importância de provas de conceito académicas como a presente.

\section{Algoritmos de Deteção e Seguimento}

No domínio da deteção e seguimento de UAVs, o estado da arte é consolidado por competições de referência internacional que estabelecem os padrões de desempenho para a comunidade científica. Em 2025, destacam-se dois eventos de particular relevância, a primeira sendo competição a WOSDETC (\emph{Workshop on Solutions for Drone Detection and Classification}), realizada no âmbito da \emph{International Joint Conference on Neural Networks} (IJCNN), cujo foco principal reside na complexa discriminação entre drones, aves em cenários de longo alcance, e o desafio Anti-UAV, coordenado pela conferência \emph{Computer Vision and Pattern Recognition} (CVPR). O desafio Anti-UAV destaca-se pela sua estrutura em três subcategorias. A primeira subcategoria foca-se na deteção e seguimento de um único drone a partir de uma posição inicial conhecida, testando a exatidão do seguidor em manter o bloqueio no alvo. A segunda subcategoria aborda a deteção global e o seguimento autónomo sem qualquer conhecimento prévio da localização do drone, exigindo que o sistema identifique o alvo no espaço aéreo de forma independente. Por fim, a terceira subcategoria foca-se no seguimento de enxames de drones, onde o desafio técnico consiste em manter a identidade única de cada unidade mesmo em situações de elevada densidade e sobreposições frequentes, garantindo que o sistema não perca o identificador individual de cada membro do enxame em cenários de sobreposição.

\cite{WOSDETC25} venceram a competição WOSDETC 2025 utilizando o modelo YOLOv11m, alcançando uma exatidão média acima de 50\% de sobreposição com o objeto real $mAP_{50}$  de 73.9\%. Para mitigar a perda de detalhe em imagens 4K na entrada nativa de 640px, dividiram os quadros em quatro segmentos sobrepostos, simulando um "efeito zoom". O \emph{dataset} foi enriquecido com a técnica \emph{copy-paste}, inserindo recortes de alvos com fundos transparentes para isolar a sua morfologia do cenário. Por fim, aplicaram interpolação linear numa janela temporal de seis quadros para recuperar deteções falhadas, garantindo um rastreio contínuo e robusto.

No desafio Anti-UAV da conferência CVPR 2025, a equipa \cite{1e2CVPR} obteve o primeiro lugar na primeira categoria e o segundo na segunda categoria. A eficácia do sistema deve-se à transformação de detetores comuns em seguidores robustos através da técnica de dinâmica de quadros (\emph{frame dynamics}), que concatena a imagem atual com mapas de diferença entre quadros. Esta abordagem permite atenuar o ruído de fundos complexos em imagens infravermelhas, destacando apenas o movimento relativo do drone. Para garantir a exatidão, o sistema utiliza uma estratégia denominada \emph{Trajectory-Constrained Filtering} (TC-Filtering), que estima a trajetória física do drone e ignora automaticamente deteções que surjam fora de um raio de movimento aceitável entre quadros. Finalmente, em vez de depender de um único algoritmo, o modelo executa um conjunto de cinco arquiteturas de alto desempenho (Cascade R-CNN, DINO, RepPoints, PAA e YOLOv11x) e aplica o método \emph{Weighted Box Fusion} (WBF). Esta técnica realiza uma fusão ponderada das \emph{bounding boxes} geradas por cada modelo, resultando numa coordenada final de alta confiança para o sistema de intersecção. Em termos métricos, o modelo alcançou uma pontuação de 73.2\% de \emph{Average Overlap Accuracy} (AOA) na primeira categoria e 57.1\% de AOA na segunda.

Na terceira categoria do desafio, focada no seguimento de múltiplos UAVs, destacou-se o trabalho de \cite{Dist-track_3CVPR} com o sistema Dist-Tracker. Este modelo utiliza a arquitetura YOLOv12 otimizada pela função de perda \emph{Scale-Shape-Quality} (SSQ) e pelo algoritmo de seguimento \emph{Fusion of L2-IoU Tracker} (FLIT). A função SSQ aborda a instabilidade geométrica típica de alvos minúsculos; nestes casos, a reduzida dimensão do drone (ex: 20 píxeis) torna o sistema extremamente sensível, onde um erro marginal de um píxel resultaria numa alteração agressivo dos pesos da rede. Para mitigar isto, o SSQ ajusta dinamicamente a função de perda com base na escala do objeto, ignora o ruído térmico de fundo e aplica penalizações específicas para UAVs com proporções horizontais ou verticais distintas (\emph{Shape-Aware}). Complementarmente, o algoritmo FLIT gere a associação de dados entre quadros. Enquanto os métodos tradicionais dependem apenas da sobreposição de caixas (IoU), o FLIT introduz uma métrica híbrida que, quando o drone se move a velocidades elevadas e não existe sobreposição direta entre quadros, utiliza o Filtro de \emph{Kalman} para prever a posição futura e associa a deteção através da menor distância euclidiana (L2) entre os centros. Esta robustez permitiu ao modelo alcançar um $AP_{50}$ de 93.9\% e uma exatidão de seguimento para alvos múltiplos (MOTA) de 77.5\%.

Além das arquiteturas CNN consolidadas, têm surgido novas abordagens fundamentadas em \emph{transformers}, a tecnologia base que permitiu o avanço dos \emph{large language models} (LLM). No domínio da deteção de objetos, destaca-se o RT-DETR (\emph{Real-Time Detection Transformer}), proposto por \cite{RT-DETR}. Esta arquitetura resolve uma das principais limitações dos \emph{transformers} anteriores de elevado custo computacional permitindo a operação em tempo real sem sacrificar a exatidão. No estudo de referência, o RT-DETR apresentou um $AP_{50}$ de 72.7\%, superando os 71.0\% obtidos pelo YOLOv8-X. A distinção fundamental reside no facto de as CNNs operarem através de kernels locais, enquanto os \emph{transformers} utilizam mecanismos de auto-atenção global. Isto permite ao modelo processar as relações entre todos os píxeis da imagem simultaneamente, capturando o contexto global e facilitando a distinção entre um drone minúsculo e o ruído visual de fundo. Embora este modelo tenha sido treinado para a deteção de objetos genéricos, é possível, tal como na família YOLO, aplicar técnicas de \emph{transfer learning} para especializar a rede na deteção exclusiva de drones.